\section{Software implementation}\label{sec:software-implementation}

To facilitate a straightforward application of our approach we have implemented both the parametric and semi-parametric likelihood ratio tests in the R package \texttt{TruncComp} which is available at the GitHub repository \citep{TruncCompGitHub}. We illustrate its applicability based on an example data set also available from the package.

The example data set can be loaded by writing \texttt{data("TruncCompExample")} after loading the package. The data set contains two variables, $Y$ and $R$, where $Y$ is the outcome and $R$ is the binary treatment indicator. There are $25$ observations in each treatment group, and truncated observations in $Y$ have been assigned the atom $\mathcal{E} = 0$. Figure~\ref{fig:dataExampleHistogram} shows histograms of the outcome for each treatment group. Visually there appears to be a difference between the two groups both in terms of the frequency of the atom at zero and a location shift in the continuous part.

\begin{lstlisting}
library(TruncComp)
data("TruncCompExample")
\end{lstlisting}

\begin{figure}[htbp]
\centering
\includegraphics[width=0.7\textwidth]{example-histogram.pdf}
\caption{Histograms of the outcome for the example data stratified by treatment group.}
\label{fig:dataExampleHistogram}
\end{figure}

The observed difference in means for the combined outcome, $\Delta$ in Equation~(\ref{eq:Delta}), is \ExampleDelta{} and both a two-sample t-test and a Wilcoxon rank sum test show highly insignificant effects of the treatment with p-values of \ExampleTTestP{} and \ExampleWilcoxP{} respectively. In order to analyse the data using proposed method we use the function \texttt{truncComp} as follows.

\begin{lstlisting}
model <- truncComp(Y ~ R, atom = 0, data = TruncCompExample, method = "SPLRT")
summary(model)
\end{lstlisting}

The output from the call to \texttt{summary} displays the estimates for the two treatment contrasts corresponding to $\mu_\delta$ and $\exp(\alpha_\delta)$ in Equations~(\ref{eq:glm1}) and~(\ref{eq:glm2}). These contrasts quantify the difference in means among the observed outcomes and the odds ratio of being observed respectively in accordance with the model specification. Each estimated treatment contrast is accompanied by a 95\% confidence interval, and finally the output displays the joint likelihood ratio test statistic and the associated $p$-value for the joint null hypothesis of no treatment effect.

From the output we see that the semi-parametric likelihood ratio analysis reports an extremely low $p$-value for null hypothesis of no joint treatment effect. This strongly contradicts the conclusions from both the previous t-test and Wilcoxon test. The confidence intervals for the two treatment contrasts indicate that the average value for the observed outcome in the group defined by $R = 1$ is significantly higher than the average value for the observed outcome in the group with $R = 0$.

The confidence intervals can also be obtained by calling the function \texttt{confint} on the model object. This command reports both the marginal confidence intervals for the two treatment contrasts as well as their simultaneous confidence region. To obtain the simultaneous confidence region we write

\begin{lstlisting}
confint(model, type = "simultaneous", plot = TRUE, offset = 1.4, resolution = 50)
\end{lstlisting}

where \texttt{resolution} is the number of discrete grid points on which the likelihood surface is evaluated for plotting. Figure~\ref{fig:simultConfidence} shows a heat-map of the semi-parametric likelihood surface as well as the 95\% simultaneous confidence region. The point $(0,0)$ is far outside of the joint confidence region which corresponds to the strong rejection of the joint null hypothesis of no treatment effect.

\begin{figure}[htbp]
\centering
\includegraphics[width=0.62\textwidth]{example-simultaneous-confidence.pdf}
\caption{Simultaneous empirical likelihood ratio surface for the two treatment contrasts.}
\label{fig:simultConfidence}
\end{figure}
